\section{Giới thiệu}
\begin{onecolentry}
    \begin{highlights}
        \item Kỹ sư Hệ thống \& Lập trình viên Backend với \textbf{hơn 3 năm kinh nghiệm phát triển phần mềm bằng Golang}, kết hợp chuyên môn sâu về hạ tầng doanh nghiệp (IDC, Ảo hóa) với kiến trúc phần mềm hiệu năng cao
        \item Hạ tầng: Máy chủ \& Lưu trữ (SAN/NAS, Dell/HPE), Ảo hóa (VMware/Proxmox, K8s), Microsoft Core (AD DS), và Bảo mật Datacenter (Firewall, WAF)
        \item Phần mềm \& DevOps: Chuyên gia Backend (Golang, Python, gRPC, DDD), CI/CD Pipeline, SRE, và Tự động hóa Hạ tầng
    \end{highlights}
\end{onecolentry}


\section{Học vấn}
\begin{twocolentry}{\textit{2019 -- 2023}}
    \textbf{Trường Đại học Bách Khoa TP.HCM - Đại học Quốc gia TP.HCM (VNU-HCM)}\\[4pt]
    \textit{Cử nhân Khoa học Máy tính}
\end{twocolentry}


\section{Ngoại ngữ}
\begin{onecolentry}
    \begin{highlights}
        \item Tiếng Anh, IELTS: 6.0
    \end{highlights}
\end{onecolentry}

\section{Kinh nghiệm Làm việc}

% --- MobiFone (Current) ---
\begin{twocolentry}{\textit{Tháng 4 2025 -- Hiện tại}}
    \textbf{Tổng công ty Viễn thông MobiFone} \\[4pt]
    \textit{Kỹ sư Hệ thống \& Lập trình viên Backend}
\end{twocolentry}
\begin{onecolentry}
    \begin{highlights}
        \item \textbf{Phần mềm \& Dữ liệu}: Phát triển Core Microservices bằng \textbf{Golang} cho luồng xử lý dữ liệu viễn thông quy mô lớn, tốc độ cao (high-throughput); duy trì \& tối ưu hệ thống xác thực Smart OTP bằng Python đảm bảo độ trễ thấp cho hàng triệu user.
        \item \textbf{Hạ tầng \& K8s}: Quản trị cụm ảo hóa (VMware, Proxmox) có tính khả dụng cao (HA, DRS) trên Server vật lý \& SAN/NAS. Quản lý container workloads và triển khai dịch vụ trên Kubernetes clusters nhằm tự động hóa failover và scale.
        \item \textbf{Mạng \& Bảo mật}: Cấu hình định tuyến mạng (VLAN, Subnetting), Active Directory (AD DS), DNS/DHCP và hệ thống an toàn thông tin (NGFW, WAF, VPN, Layer 7 Load Balancers).
        \item \textbf{SRE \& Giám sát}: Triển khai giám sát tập trung (LibreNMS, Zabbix, Prometheus, Grafana) để tự động cảnh báo sự cố, giảm thiểu MTTR và đảm bảo SLA. Quản trị chiến lược Disaster Recovery (Veeam, SRM) tuân thủ RPO/RTO.
        \item \textbf{CI/CD \& Tự động hóa}: Xây dựng các luồng triển khai tự động (CI/CD pipelines) và áp dụng tư duy Infrastructure as Code (IaC) với Ansible để đạt được zero-downtime deployments.
    \end{highlights}
\end{onecolentry}
\vspace{0.2 cm}

% --- ActsOne ---
\begin{twocolentry}{\textit{Tháng 10 2023 -- Tháng 3 2025}}
    \textbf{ActsOne} \\[4pt]
    \textit{Kiến trúc sư Giải pháp \& Kỹ sư Hệ thống}
\end{twocolentry}
\begin{onecolentry}
    \begin{highlights}
        \item \textbf{Nền tảng Xử lý Đơn hàng (Backend):} Xây dựng hệ thống backend (Golang, gRPC, DDD) tích hợp API Shopee, Lazada, TikTok với cơ chế cron/webhook, Redis caching và khóa phân tán (distributed locking) để xử lý lượng đơn hàng khổng lồ theo thời gian thực.
        \item \textbf{Hạ tầng Ảo hóa IDC (System):} Thiết kế và vận hành toàn bộ hạ tầng on-premise với cụm Proxmox VE. Ứng dụng Ansible để tự động hóa cấp phát VM (nhanh hơn 80\%), thiết lập cô lập VM, snapshot và tự động failover (HA).
        \item \textbf{Nền tảng Kubernetes \& Lưu trữ (DevOps/System):} Tự động hóa triển khai các cụm K8s/K3s bằng Ansible playbooks. Cấu hình lưu trữ bền vững (persistent storage) với Longhorn và chiến lược lưu trữ phân tầng (hot/warm/cold archive) trên NAS, giúp giảm 40\% chi phí lưu trữ.
        \item \textbf{Cơ sở Dữ liệu \& SRE:} Thiết lập MySQL/PostgreSQL replication (master-slave, multi-master) và Redis caching, tăng 50\% thông lượng hệ thống. Xây dựng nền tảng SRE giám sát toàn diện, thiết lập sao lưu tự động (Backup \& Restore) cho VM/DB để đạt SLA 99.9\%.
        \item \textbf{Mạng \& Bảo mật Zero-Trust:} Kiến trúc mạng an toàn toàn diện tích hợp DrayTek firewalls, Cloudflare ZeroTrust (WAF, chống DDoS), NGINX load balancer và mạng riêng ảo VPN (Headscale, Tailscale/WireGuard).
        \item \textbf{Identity Provider \& API Gateway:} Triển khai Authentik làm IdP trung tâm (SSO, MFA, RBAC). Phát triển hệ thống API Gateway bằng Golang với rate limiting (50 RPM/client) và document Swagger/Scalar, giảm 85\% băng thông API rác.
        \item \textbf{CI/CD Pipelines (DevOps):} Xây dựng hoàn chỉnh các luồng CI/CD (GitHub Actions, Drone CI) để tự động hóa quá trình kiểm thử, build và deploy microservices lên Kubernetes.
    \end{highlights}
\end{onecolentry}
\vspace{0.2 cm}

% --- Kays.vn ---
\begin{twocolentry}{\textit{2023 (Làm việc từ xa)}}
    \textbf{Kays.vn} \\[4pt]
    \textit{Kỹ sư Full-stack \& Backend}
\end{twocolentry}
\begin{onecolentry}
    \begin{highlights}
        \item Xây dựng các module backend đa nền tảng bằng \textbf{Rust} để xử lý giao dịch, hỗ trợ mã hóa bảo mật cao cấp (Tauri Dashboard).
        \item Triển khai luồng dữ liệu (data pipeline) theo mô hình hướng sự kiện (event-driven) bằng AWS EventBridge để đồng bộ dữ liệu thời gian thực giữa Supabase và hệ thống cũ.
        \item Thiết kế và phát triển RESTful APIs bằng \textbf{Golang} cho phân hệ sản phẩm và đơn hàng. Phát triển frontend bằng Next.js giúp tối ưu hóa SEO.
        \item Đào tạo (mentor) 4 thực tập sinh về các best practice trong TypeScript, Golang và kỹ thuật tối ưu hóa hiệu năng.
    \end{highlights}
\end{onecolentry}
\vspace{0.2 cm}

% --- IVS ---
\begin{twocolentry}{\textit{Tháng 8 2022 -- Tháng 9 2023}}
    \textbf{IVS Technology Solutions} \\[4pt]
    \textit{Lập trình viên Web \& Backend}
\end{twocolentry}
\begin{onecolentry}
    \begin{highlights}
        \item Phát triển các RESTful APIs cốt lõi bằng SpringBoot (Java) và xây dựng giao diện động bằng Next.js/React.js (mobile-first design).
        \item Tối ưu hóa truy vấn cơ sở dữ liệu MongoDB/MySQL thông qua indexing, giảm 15\% thời gian tải trang.
        \item Ứng dụng kỹ thuật Server-Side Rendering (SSR) và lazy loading để phân phối nội dung tốc độ cao.
        \item Xây dựng quy trình CI/CD bằng Jenkins để tự động hóa toàn bộ việc kiểm thử và triển khai (deployment).
    \end{highlights}
\end{onecolentry}

\section{Kỹ năng Công nghệ \& Lĩnh vực Chuyên môn}
\begin{onecolentry}
\begin{highlights}
    \item \textbf{Lập trình:} Go (Golang), C/C++, Python, JavaScript, Java, Rust
    \item \textbf{Backend \& API:} gRPC, REST, Swagger, Scalar, Domain-Driven Design (DDD), Microservices, API Gateway
    \item \textbf{Hạ tầng (System):} On-Prem IDC, Proxmox VE, Kubernetes (K8s/K3s), Colocation, Platform Engineering
    \item \textbf{Kiến trúc \& SRE:} End-to-End System Design, HA Architecture, Monitoring, SLA Management
    \item \textbf{Cơ sở dữ liệu:} MySQL/PostgreSQL (Replication, Master-Slave, Multi-Master), MongoDB, Redis, DuckDB
    \item \textbf{Lưu trữ (Storage):} Hệ thống NAS, Hot/Warm/Cold Archive, Storage Tiering, Data Lifecycle Management
    \item \textbf{Bảo mật \& Định danh:} Cloudflare ZeroTrust, WAF, Secure Tunnels, DrayTek Firewalls, Authentik (SSO, MFA, RBAC)
    \item \textbf{Mạng (Networking):} VLAN, Load Balancing, Network Segmentation, VPN (Tailscale/WireGuard)
    \item \textbf{DevOps \& CI/CD:} Docker, Ansible, GitHub Actions, Drone CI, Jenkins, Nginx, Infrastructure as Code (IaC)
    \item \textbf{Web UI/Client:} Next.js, React.js, SpringBoot, WordPress, Tauri, Electron
\end{highlights}
\end{onecolentry}

\section{Thành tựu \& Giải thưởng}
\begin{onecolentry}
    \begin{highlights}
        \item \textbf{Giải pháp Tốt nhất (Best Solution Award)} - ActsOne (2024)
        \item \textbf{Nhân viên Xuất sắc (Outstanding Employee)} - Kays.vn (2023)
        \item \textbf{Thực tập sinh Triển vọng (Most Promising Intern Award)} - IVS Technology Solutions (2022)
        \item \textbf{Top 36 Cuộc thi Thuyết trình} - Cấp Đại học (2019)
        \item \textbf{Học bổng Khuyến khích Học tập} - Đại học Bách Khoa TP.HCM (2019, 2022)
        \item \textbf{Giải thưởng Dự án Sinh viên (Điểm: 9.30/10)} - Đại học Bách Khoa TP.HCM (2022)
    \end{highlights}
\end{onecolentry}
